\section{机械波}
\sourceproblem{选择1}{17}{23}
\problemtext 给定行波 $y=6\cos(\pi t-3\pi x+\pi/2)$（SI），判断波速与频率。
\workings $\omega=\pi,\ k=3\pi$，所以 $u=\omega/k=1/3\,\mathrm{m/s}$，$f=\omega/(2\pi)=1/2\,\mathrm{Hz}$，沿正 $x$ 传播，选B。
\crosscheck 波速是相位传播速度，不是质点横向振动速度。
\sourceproblem{选择2}{17}{23}
\problemtext 图1为正向传播的瞬时波形，判断波长及 $x=6\,\mathrm m$ 处质点运动方向。
\workings 图中相邻波峰相距8米，$x=6\,\mathrm m$ 处波形斜率为正。正向行波满足 $\partial_t y=-u\partial_x y$，故该处质点向下运动，选D。
\crosscheck 单张波形图在不知道波速时不能单独给出周期。
\sourceproblem{选择3}{17}{23}
\problemtext 图1的波速为10米每秒，求频率。
\workings $\lambda=8\,\mathrm m$，$f=u/\lambda=1.25\,\mathrm{Hz}$，选A。
\crosscheck $T=\lambda/u=0.8\,\mathrm s$，不是0.8赫兹。
\sourceproblem{选择4}{17}{23}
\problemtext 正向行波周期0.5秒、振幅1米、波长2米，原点初始处于负最大位移。
\workings $\omega=4\pi,k=\pi,\varphi_0=\pi$，故 $y=\cos(4\pi t-\pi x+\pi)$（SI）。
\crosscheck 原点初位移为负1米，速度为零；固定相位随时间向右移动。
\sourceproblem{选择5}{17}{23}
\problemtext 图2为沿负轴传播的波在 $t=T/4$ 时的波形，求原点振动初相。
\workings 图中原点位移为零、空间斜率为正。左行波满足 $\partial_t y=u\partial_x y>0$，故该时刻余弦相位为 $3\pi/2$，减去 $\omega T/4=\pi/2$ 得初相 $\pi$，选C。
\crosscheck 采用余弦初相约定；原点在初始时处于负最大位移。
\sourceproblem{选择6}{17}{23}
\problemtext 图3给定 $t=T/4$ 时的波形，求P点的振动方程。
\workings 图中振幅2厘米，右行波速20米每秒，波长按图刻度读为4米，则 $\omega=10\pi$。P处为正斜率过零点，质点向下运动；此刻余弦相位 $\pi/2$，减去 $\pi/2$ 得初相零。因此 $y_P=2\cos(10\pi t)$（cm）。
\crosscheck 若读图的P点或刻度与选项不符，应按 $y_P(t)=A\cos[\omega(t-T/4)+\phi_P]$ 重新辨识，$\cos\phi_P=y_P(T/4)/A$ 且速度符号由右行波斜率确定。本题保留读图待核标记，不强配选项。
\sourceproblem{选择7}{17}{23}
\problemtext 正向行波波长4米，图4中A至B相距5米，A的振动为 $5\cos(2\pi t+\pi)$（SI）；求B的振动。
\workings 相位滞后 $2\pi(5/4)=5\pi/2$，故 $y_B=5\cos(2\pi t+\pi/2)$（SI），选A。
\crosscheck 传播时间为 $5/u$，相位差与相应时间延迟一致；图示在PDF24。
\sourceproblem{选择8}{18}{24}
\problemtext 左行波波长 $\lambda$，比较 $x_1=\lambda,x_2=-\lambda/2$ 两点相位。
\workings 同时相位为 $\omega t+kx+\varphi$，所以 $\phi_1-\phi_2=k(3\lambda/2)=3\pi$；若反向定义相位差则为 $-3\pi$。两点反相。
\crosscheck 必须说明相位差的减法顺序；模 $2\pi$ 意义下等效为 $\pi$。
\sourceproblem{选择9}{18}{24}
\problemtext 理想弹性介质中的平面简谐行波，判断质元的动能、势能关系。
\workings 在无耗散线性一维行波模型，$w_K=\rho(\partial_t y)^2/2$，$w_U=\rho u^2(\partial_x y)^2/2$。因 $\partial_t y=\pm u\partial_x y$，二者始终相等，选D。
\crosscheck 质元在位移极值处两种能量密度同时为零，在过平衡位置时同时最大；不同于孤立弹簧振子。
\sourceproblem{选择10}{18}{24}
\problemtext 两同相相干波源相距 $\lambda/2$，图5直线外侧已知相消，判断连线上干涉。
\workings 两源外侧任一点的路程差恒为 $\lambda/2$，相位差为奇数倍 $\pi$，为相消条件；两源之间路程差连续从 $-\lambda/2$ 变到 $\lambda/2$，中点相长，不能说内部各点均相消或均相长。
\crosscheck 完全抵消还需两波在该点振幅相等；只根据相位判断增强或减弱。
\sourceproblem{选择11}{18}{24}
\problemtext 判断驻波的能量和相位命题。
\workings 理想驻波可写成 $y=2A\sin kx\cos\omega t$，相邻两波节之间 $\sin kx$ 同号，故非节点质元同相；跨过波节相位反转。选C。
\crosscheck 波节位移恒零，但介质形变能不必为零；波腹位移最大也不代表任意时刻总能量密度最大。
\sourceproblem{填空1}{18}{24}
\problemtext 行波 $y=5\cos(\pi t+4\pi x-\pi/2)$（SI），求传播方向、振幅、周期、波长和波速。
\workings 负 $x$ 方向，$A=5\,\mathrm m$，$T=2\,\mathrm s$，$\lambda=1/2\,\mathrm m$，$u=1/4\,\mathrm{m/s}$。
\crosscheck 固定相位满足 $\pi+4\pi\,\mathrm dx/\mathrm dt=0$。
\sourceproblem{填空2}{18}{24}
\problemtext 波在介质A中速率为 $v$，进入B后波长为原波长两倍，求B中波速。
\workings 静止界面两侧频率相同，$u_B=f(2\lambda_A)=2v$。
\crosscheck 波速改变由介质决定；频率由波源决定。
\sourceproblem{填空3}{18}{24}
\problemtext 图6瞬时波形给出A点向上运动，判断传播方向以及B、C、D点振动方向。
\workings A处空间斜率为负，由 $\partial_t y=-u\partial_x y$ 得波沿正 $x$ 传播。B处为正斜率，向下；C为负斜率，向上；D在波峰，瞬时速度为零，随后向下。
\crosscheck 原图D的精确位置待独立读图复核；若实际标在下降段而非峰顶，则D瞬时速度向上。严格判据为局部空间斜率而非位移正负。
\sourceproblem{填空4}{18}{24}
\problemtext $y=7\cos(3t-5x)$（SI），求 $x=-3$ 处振动式及与 $x=4$ 处的相位差。
\workings $y_{-3}=7\cos(3t+15)$，$y_4=7\cos(3t-20)$，所以 $\phi_{-3}-\phi_4=35\,\mathrm{rad}$。
\crosscheck 若要求主值，可对35弧度取模 $2\pi$；相位差符号须声明顺序。
\sourceproblem{填空5}{18}{24}
\problemtext 图7两波源初相差 $\pi$，到P距离分别为 $a,b$（$b>a$），P处相长干涉，求可取波长。
\workings 到P相位差为 $\pi\pm 2\pi(b-a)/\lambda$，相长要求为 $2n\pi$，故
\[
\lambda=\frac{2(b-a)}{2j+1},\qquad j=0,1,2,\ldots .
\]
\crosscheck 初相反相时奇数个半波长的路程差使到达相位一致；若问最长波长则取 $j=0$。
\sourceproblem{填空6}{18}{24}
\problemtext 同频平面简谐波强度之比为 $a$，求振幅之比。
\workings 在同一线性介质、同波型下 $\bar I=\rho u\omega^2A^2/2$，故 $A_1/A_2=\sqrt a$。
\crosscheck 若两介质波阻抗不同，需另乘波阻抗比的平方根，不能仅由同频推出该结果。
\sourceproblem{计算1}{19}{25}
\problemtext 左行简谐波波长6米，图8给出 $x=3\,\mathrm m$ 处振动曲线，求该点、行波及原点振动式。
\workings 由图读振幅10厘米、相邻峰值时间差约12秒，$t=0$ 处为正位移且下降；设从图确认的初相为 $\phi_3$。则
\[
y_3=10\cos(\omega t+\phi_3),\quad
y(x,t)=10\cos[\omega t+\tfrac{\pi}{3}(x-3)+\phi_3],\quad
y_0=10\cos(\omega t+\phi_3-\pi)
\]
（cm）。若峰值在 $t=10\,\mathrm s$ 且周期12秒，则 $\omega=\pi/6,\phi_3=\pi/3$（模 $2\pi$）。
\crosscheck 原图时间刻度较细，周期及初相保留条件式待复核；波长6米决定空间相位系数 $\pi/3$，不受读图时间误差影响。
\sourceproblem{计算2}{19}{25}
\problemtext 正向行波波速 $u$，图9中A位于 $x=a$、B位于 $x=3a$，A处振动 $y_A=A_0\cos\omega t$。求O、B处振动及与B同状态点。
\workings $k=\omega/u$，
\[
y(x,t)=A_0\cos[\omega t-k(x-a)],\quad
y_O=A_0\cos(\omega t+ka),\quad
y_B=A_0\cos(\omega t-2ka).
\]
与B同相且振幅相同的位置为 $x=3a+n(2\pi u/\omega)$，$n\in\mathbb Z$。
\crosscheck 同状态包括位移和速度，要求相位差为整周；若波尚未传播到某点，该点不能计入已振动区域。

\sourceproblem{计算3}{20}{26}
\problemtext 图10给定正向简谐行波初始波形，波速10米每秒，求原点振动及波动方程。
\workings 放大图10可见振幅10厘米，原点与 $x=10\,\mathrm m$ 处为同向空间过零点，故 $\lambda=10\,\mathrm m$，$k=\pi/5\,\mathrm{rad/m}$，$\omega=ku=2\pi\,\mathrm{rad/s}$。初始波形为 $y(x,0)=0.10\sin(\pi x/5)$，沿正向平移得
\[
y(x,t)=0.10\sin(\pi x/5-2\pi t),\qquad
y(0,t)=-0.10\sin(2\pi t)
\]
（SI）。
\crosscheck 原点初始空间斜率为正，故初速度为负；$y(x,t)=y(x-ut,0)$ 与传播方向一致。
\sourceproblem{计算4}{20}{26}
\problemtext 正向行波波速2米每秒、周期0.2秒、振幅3厘米，初始 $x=1\,\mathrm m$ 处位移为负最大值，求波动方程。
\workings $\lambda=uT=0.4\,\mathrm m$，$\omega=10\pi$，$k=5\pi$。令 $y=0.03\cos(10\pi t-5\pi x+\varphi)$，由 $-5\pi+\varphi\equiv\pi\pmod{2\pi}$ 得 $\varphi\equiv0$，所以
\[
y(x,t)=0.03\cos(10\pi t-5\pi x)\quad(\mathrm{SI}).
\]
\crosscheck 代入 $x=1,t=0$ 得 $-0.03\,\mathrm m$，此时质点速度为零；相位传播速率为2米每秒。
